\documentclass[10pt]{beamer}

\usepackage{verbatim}
%\usepackage[latin1]{inputenc}
\usepackage{amssymb,amsmath,amsfonts,url}
\usepackage{graphicx,color}
\usepackage{pgf,pgfarrows,pgfnodes,pgfautomata,pgfheaps,pgfshade}
\usepackage[utf8]{inputenc}
\usepackage[french]{babel}
\usepackage{tikz}

\usepackage[all]{xy}
\usepackage{amstext}
\usepackage{xspace}
\usepackage{mdwtab}


\mode<article> {
  \usepackage{fullpage}
  \usepackage{hyperref}
  \setjobnamebeamerversion{advancedcourse.beamer}
}





\newcommand{\alertonce}[1]{\addtocounter{beamerpauses}{-1}\alert<+>{#1}}
\newcommand{\alerttwice}[1]{\addtocounter{beamerpauses}{-1}\alert<+-+(1)>{#1}}

%\newcommand{\subsubsubsection}[1]{\noindent\textbf{#1}}

\newcommand{\cfbox}[1]{\begin{center}\fbox{#1}\end{center}}
\newcommand{\cbox}[1]{\begin{center} #1 \end{center}}

%%%%%%%%%%%%%%%%%%%%%%%%%%%%%%%%%%%%%%%%%%%%%%%%%%%%%%%%%%%%%%%%%%%%%%%%%%%%%%%%%%%%%%%%%

\graphicspath{{Graphics/}}
%\input def_0.tex

%\input{StyleSlides.tex}
\mode<presentation> {
    \usepackage{beamerthemeBoadilla}

  %\beamertemplatetransparentcovereddynamic
  \beamertemplateballitem
}


\title[IN100 Cours 7]{\textbf{Fondements informatiques I}\\
Cours 7: Types et typage}

\author[]{Sorina Ionica \texttt{sorina.ionica@uvsq.fr} \\~\\Sandrine Vial  \texttt{sandrine.vial@uvsq.fr} }


\institute{}

\date{}


%\setbeamercolor{normal text}{fg=black}

\begin{document}


\frame{\titlepage}


\frame[containsverbatim]{
\frametitle{Les séquences}

\textcolor{blue}{\`A retenir} : Une séquence est ensemble fini et ordonné d'éléments indicés de $0$ à $n-1$ (si la séquence comporte $n$ éléments).

\begin{itemize}
\item les listes : modifiables
\item les tuples:  non-modifiables
\item les chaînes de caractères : non-modifiables
\item les "range" (intervalles)
\end{itemize}

\vspace{0.5 cm}

\small{
\begin{center}
\begin{tabular}{ |c|c| } 
 \hline
\texttt{x in s} &  \texttt{True} si un élément de \texttt{s} est égal à \texttt{x}, sinon \texttt{False}  \\ 
\texttt{x not in s} & \texttt{False} si un élément de \texttt{s} est égal à \texttt{x}, sinon \texttt{True} \\ 
\texttt{s + t}  &  la concaténation de \texttt{s} et \texttt{t}  \\ 
\texttt{s * n} ou \texttt{n * s} & ajouter \texttt{s} $n$ fois à lui-même \\
\texttt{s[i]} & $i$-ème élément de \texttt{s} en commençant par 0 \\
\texttt{s[i:j]} & tranche (slice) de \texttt{s} de \texttt{i} à \texttt{j-1} \\
\texttt{s[i:j:k]} & tranche (slice) de \texttt{s} de \texttt{i} à \texttt{j-1} avec un pas de \texttt{k} \\
\texttt{len(s)} &  longueur de \texttt{s} \\
 \hline
\end{tabular}
\end{center}}


}

\frame[containsverbatim]{
\frametitle{Les séquences}

\textcolor{blue}{\`A retenir} : Les séquences sont des objets \textbf{itérables.}\\


 On peut  boucler sur tous éléments grâce à une boucle \textbf{for}.



\begin{exampleblock}{}
\small{
\begin{verbatim}
ma_chaine = "Abracadabra"

for carac in ma_chaine:
  print(carac)
  
l = [55, 23,103,258]
t = (128, 0, 255)

for nb in l:
    print(nb)

for nb in t:
    print(nb)
 \end{verbatim}
 }
\end{exampleblock}


}


\begin{frame}[fragile]{Chaînes de caractères}


\begin{itemize}
%\item  Les chaînes de caractères peuvent être vues comme une sorte de listes particulières.
\item 
On peut notamment utiliser  les \emph{tranches}.
\end{itemize}

\begin{exampleblock}{}
\small{
\begin{verbatim}
chaine = 'Abracadabra!'
print(chaine[2:5])
print(chaine[-1])
\end{verbatim}
}
\end{exampleblock}

\vspace{1 cm}

\begin{itemize}
\item  La différence principale entre les listes et les chaînes de caractères
est que les chaînes de caractères sont des objets
\textbf{non-modifiables} après création.
\end{itemize}

\begin{exampleblock}{}
\small{
\begin{verbatim}
chaine[2] = 'd'  #TypeError: 'str' object does not support item assignment 
liste_chaines=["Paris","londres", "Rome","Budapest"]
liste_chaines[1][0]='L'
\end{verbatim}
}
\end{exampleblock}

\end{frame}

\begin{frame}[fragile]{Les encodages}

\small{En Python, un caractère est une chaîne de longueur 1. \\}

\begin{block}{\small{L'encodage \textbf{ASCII} (American Standard Code for Information Interchange)}}
\small{
 ASCII représente 128 caractères anglais sous forme de nombres, chaque lettre étant associée à un numéro spécifique compris entre 0 et 127. 
 \begin{itemize}
 \item Par exemple, le code ASCII pour la lettre majuscule A est 65, pour la majuscule B c’est 66 etc. 
 \end{itemize}
 }
\end{block}    
    

\small{
\begin{itemize}
\item La fonction \texttt{ord} convertit un caractère dans sa représentation numérique. 
\item La fonction \texttt{chr} convertit un nombre dans un caractère ASCII. 
\end{itemize}
}

\pause
\vspace{0.5 cm}
\small{
En fait, Python repose sur l'encodage Unicode, qui comprend beaucoup plus de caractères: 
\begin{exampleblock}{}
\begin{verbatim}
>>> chr(245)
'õ'
\end{verbatim}
\end{exampleblock}
}
\end{frame}


\begin{frame}[fragile]{Les ensembles}

Un \textbf{ensemble (set)} est une collection d'éléments ne contenant
pas de doublons.

\vspace{0.3 cm}

\textbf{Exemple:}  $A = \{1, 3, 10, 5, 2\}$ est un ensemble, tandis que
$B = \{1, 10, 3, 3, 5, 2, 1\}$ n'en est pas un.

\vspace{0.3 cm}

Python permet de créer des objets de type \texttt{set}.

\begin{itemize}
\item Un \texttt{set} n'est ni \textbf{ordonné}, ni \textbf{modifiable}.
\item On utilise des \texttt{\{\}} pour créer un nouveau \texttt{set}.
\end{itemize}

\vspace{0.3 cm}
\begin{exampleblock}{}
\small{
\begin{verbatim}
s = {1, 2, 3, 4, 5}
print(s)
print(type(s))
\end{verbatim}
}
\end{exampleblock}
\end{frame}




\frame[containsverbatim]{
\frametitle{Créer des ensembles}


La fonction \texttt{set()} permet de générer un ensemble à partir de
n'importe quel objet itérable.

\vspace{0.5 cm}

\begin{exampleblock}{}
\begin{verbatim}
# À partir d'une liste
s = set([1, 3, 4, 5, 4, 3])

# À partir de range()
s = set(range(10))

# À partir d'une chaîne de caractères
s = set("Versailles")
\end{verbatim}
\end{exampleblock}

}


\frame[containsverbatim]{
\frametitle{Accéder aux éléments d'un ensemble}

\begin{itemize}
\item  Un \texttt{set} est \textbf{itérable}. 
 \item Un set n'est pas \textbf{ordonné}, donc il n'est pas une séquence. 
 \begin{itemize}
\item  Il est impossible de récupérer un élément par sa position.
\end{itemize}
\end{itemize}

\begin{exampleblock}{}
\begin{verbatim}
s = {0, 5, 10, 15}
for element in s :
    print(element)
    
print(s[2])     #TypeError: 'set' object is not subscriptable
\end{verbatim}
\end{exampleblock}


}
  
  
\begin{frame}[fragile]{Doublons}  


Les \texttt{sets} sont très pratiques pour \emph{supprimer des
doublons}.

\vspace{0.5 cm}

\textbf{Exercice} : Écrire une fonction qui prend en paramètre une liste
et renvoie une nouvelle liste, égale à celle passée en paramètre mais
sans les doublons.

\pause

\vspace{0.5 cm}

\begin{exampleblock}{}
\begin{verbatim}
def supprime_doublons(liste) :
    return list(set(liste))

nouvelle_liste = supprime_doublons([1, 3, 3, 5, 6, 7, 3, 1])
print(nouvelle_liste)
\end{verbatim}
\end{exampleblock}


\end{frame}




\frame[containsverbatim]{
\frametitle{Opérateurs sur les ensembles}


Les ensembles permettent d'effectuer des opérations mathématiques telles
que l'\emph{union}, l'\emph{intersection}, la \emph{différence} ou la
\emph{différence symétrique}.

\vspace{0.5 cm}

\textbf{Exemple:} \(A = \{0, 1, 2, 3, 4\}\), \(B = \{2, 3, 4, 5\}\)

\begin{itemize}
\item
  \(A\cup B := \{x \in A \text{ ou } x\in B\} = \{0, 1, 2, 3, 4, 5\}\)
\item
  \(A\cap B := \{x \in A \text{ et } x\in B\} = \{2, 3, 4\}\)
\item
  \(A\setminus B := \{x \in A \text{ et } x \not \in B\} = \{0, 1\}\)
\item
  \(A\triangle B := \{x \in A\cup B \text{ et } x \not \in A\cap B\} = \{0, 1, 5\}\)
\end{itemize}

}


\frame[containsverbatim]{
\frametitle{Opérateurs sur les ensembles}



 \begin{itemize}
\item
  \texttt{a\ \textbar{}\ b}: \(A\cup B\)
\item
  \texttt{a\ \&\ b}: \(A\cap B\)
\item
  \texttt{a\ -\ b}: \(A \setminus B\)
\item
  \texttt{a\ \^{}\ b}: \(A\triangle B\)
\end{itemize}


\begin{exampleblock}{}
\begin{verbatim}
chaine1 = 'chaussette'
chaine2 = 'archiduchesse'

a = set(chaine1)
b = set(chaine2)

print(a | b)
print(a & b)
print(a - b)
print(a ^ b)
\end{verbatim}
\end{exampleblock}




}

\frame[containsverbatim]{
\frametitle{Dictionnaires}


Un dictionnaire est une structure de données, dont les éléments sont
accessibles grâce à une \emph{clé}.

\vspace{0.5 cm}

 \begin{itemize}
\item Pour construire un dictionnaire, on utilise des accolades \texttt{\{\}}.
\item Un élément d'un dictionnaire est un couple \texttt{cle:\ valeur}.
\item On accède à une \textbf{valeur} du dictionnaire en utilisant la \textbf{clé}. 
\end{itemize}


\vspace{0.5 cm}
\begin{exampleblock}{Exemple}
\small{
\begin{verbatim}
dico_notes = {"Gaëtan" : 14, "Laure" : 9, "Kenza" : 17, "Gabriel" : 14}

print(dico_notes["Laure"])
\end{verbatim}
}
\end{exampleblock}

}


\frame[containsverbatim]{
\frametitle{Dictionnaires}



\begin{itemize}
\item Dans un dictionnaire, les clés doivent être \emph{uniques}, mais pas
  les valeurs.
\item Les clés peuvent être n'importe quel objet non-modifiable (immutable) 
\begin{itemize}
\item Par example des \texttt{int} ou des \texttt{string}. 
\end{itemize}
\item Il est facile de créer et modifier un dictionnaire.
\end{itemize}



\vspace{0.5 cm}

\begin{exampleblock}{}
\small{
\begin{verbatim}
dico_notes = {"Gaëtan" : 14, "Laure" : 9, "Kenza" : 17, "Gabriel" : 14}

# Changer une entrée existante
dico_notes["Gaëtan"] = 12

# Ajouter une nouvelle entrée

dico_notes["Anne"] = 15
print(dico_notes)

# Créer un dictionnaire vide.
dico_vide = {}
\end{verbatim}
}
\end{exampleblock}


}



%# Utilisation du constructeur dict()
%dico_restos = dict([("L'auberge du Goëllo", 22170), ("Café Menilmontant", 75011), ("Mamie Bigoude", 17000)])
%print(dico_restos)


\frame[containsverbatim]{
\frametitle{Parcourir un dictionnaire}

On peut vérifier facilement qu'une valeur est dans le dictionnaire si on connaît sa clé. 


\begin{exampleblock}{}
\begin{verbatim}
cle="Kenza"
if cle in dico_notes:
     print(cle, "a obtenu la note", dico_notes[cle])     
\end{verbatim}
\end{exampleblock}


\vspace{0.5 cm}

Les clés et leurs valeurs peuvent être récupérées en même temps en
  utilisant la méthode \texttt{items()}. 

\vspace{0.5 cm}

\begin{exampleblock}{}
\small{
\begin{verbatim}
dico_notes = {"Gaëtan" : 14, "Laure" : 9, "Kenza" : 17, "Gabriel" : 14}

for i in dico_notes.items():
    print(i)
\end{verbatim}
}
\end{exampleblock}


}


\begin{frame}[fragile]{Typage en Python}


\textbf{Rappel} : En Python le typage est dynamique. 

\begin{exampleblock}{}
\begin{verbatim}
a=5
b=7
print(a+b)  # int + int
x="salut"
y="hello"
print("salut" + "hello") # str + str
print(a + y)            #int + str   Erreur!

# Problème de type pour l'opération /
"hello" / 2
\end{verbatim}
\end{exampleblock}
\end{frame}


\begin{frame}[fragile]{Annotation de type}



\begin{itemize}
\item Depuis la version 3.5 on peut ajouter \textbf{des annotations de type} aux variables et fonctions. 
\item Ces annotations  permettent d'apporter plus de structure aux programmes pour éviter certains bugs. 
\item Ces annotations sont facultatives et vous pouvez mélanger du code annoté ou non. 
\end{itemize}

\vspace{0.5 cm}


\begin{exampleblock}{}
\begin{verbatim}
a: int = 5
b: int =7
print(a+b)  # int + int
x: str="salut"
y: str = "hello"

y=6
print(type(y))

liste: list[int] = [1, 2, 3]
\end{verbatim}
\end{exampleblock}

\end{frame}


\begin{frame}[fragile]{Annotation de type pour les fonctions}



On peut spécifier le type des arguments. 
 \begin{itemize}
 \item  Les annotations \texttt{a: int}, \texttt{b:int} indiquent que les paramètres doivent être des entiers. 
 \end{itemize}

\begin{exampleblock}{Exemple}
\small{
\begin{verbatim}
def mon_print(a : int, b: int):     
	       print(a + b)
      
mon_print(5,7)
mon_print("salut","hello")
mon_print(5,"salut")
\end{verbatim}
}
\end{exampleblock}


\begin{itemize}
\item Les informations de type ne changent rien à l'exécution. 
\item  Il y a une erreur de type seulement  quand on applique une fonction à un élément et qu'il n'est pas du type attendu par la fonction.
\end{itemize}

\end{frame}


\begin{frame}[fragile]{Annotation de type}



On peut spécifier le type de retour. 
\begin{itemize}   
\item    L'annotation  \texttt{-> int} indique que la fonction renvoie un entier.
 \end{itemize}


\begin{exampleblock}{Exemple}
\small{
\begin{verbatim}
def somme(a: int, b: int) -> int:
	      return a + b


print(somme(5,7))
print(somme("salut","hello"))	
\end{verbatim}
}
\end{exampleblock}

\begin{itemize}
\item Pour détecter les erreurs, il faut utiliser des outils comme \texttt{mypy}. 
\end{itemize}



\end{frame}

\end{document}

